\chapter{The generative setting}\label{ch:grammars}

This chapter fixes the objects of the monograph: generative systems and their histories, the UNNS operator SELECT-II as the principal realisation, symbolic event grammars as a second realisation, and the permutohedral complexes that carry the topology. It ends with a contract stating which constructions need integer divisibility and which hold for arbitrary finite sets of histories.

\section{Generative systems, histories and fibres}
@@def:grammar@@

A fibre is a disjoint union of \emph{pieces} $P(x_0,Q)$ of histories with a given source and content. We write $\Ent(\cdot)$ for Shannon entropy in bits, and use the uniform distribution on a finite fibre unless another is stated.

@@def:swap@@

@@lem:swap-invariants@@

\section{The UNNS operator SELECT-II}\label{sec:selectII}
SELECT-II is the operator of the Substrate whose reachability and complete-history structure was established in the Foundations of the Substrate~\cite{UNNSFound}. Its histories can be enumerated exactly and its admissibility is governed by divisibility, which makes it a rich source of finite combinatorial structure. It is the principal worked realisation of this monograph; the monograph does not claim that it is the only generative rule of the Substrate.

@@def:selectII@@
SELECT-II is a commutative realization with $r_N=R(N)$, so \cref{lem:swap-invariants} applies. A label word is \emph{formally} applicable from any seed, since $R(N)$ is rational, but \emph{admissible} only if every intermediate state is an integer divisible by the next label.

\section{Symbolic event grammars}\label{sec:symbolic}
For the topological chapters no integer state is needed. Take $k$ distinguishable events and a set $A\subseteq S_k$ of admissible orderings. Two families occur. A \emph{pairwise grammar} is given by a digraph $\Gamma$ on the events, with $A$ the set of orderings whose consecutive pairs are arcs of $\Gamma$ (the Hamiltonian paths of $\Gamma$); when events carry categories, $\Gamma$ is the preimage of a rule set on categories. A \emph{triple grammar} forbids a set of consecutive triples. Neither has a seed or arithmetic, yet each yields a set of orderings, hence a subcomplex of the permutohedron. They are used in \cref{ch:locality} to separate what depends on arithmetic from what depends only on the \emph{locality} of the constraints. A reading with categories such as Sunny, Rainy, Snowy is a symbolic test system only; no weather data enter.

\section{Admissible sets and the permutohedral complex}\label{sec:complex}
Fix $k\ge2$, write $[k]=\{0,\dots,k-1\}$ and let $S_k$ be the set of orderings $w=(w_1,\dots,w_k)$ of $[k]$. The \emph{prefix sets} of $w$ are $I_j(w)=\{w_1,\dots,w_j\}$; the chain $I_1(w)\subset\dots\subset I_{k-1}(w)$ is the \emph{flag} of $w$, and $w\mapsto$ flag is a bijection between orderings and maximal chains of the poset $\Pm$ of proper nonempty subsets of $[k]$.

@@def:faces@@
The $1$-skeleton of $P_k$ is the Cayley graph of $S_k$ for the adjacent transpositions $s_1,\dots,s_{k-1}$ (positions $j,j{+}1$); an edge is a face with one block of size two, and the graph is bipartite by the sign of the permutation. We use the standard facts about permutohedra and Coxeter groups~\cite{Ziegler,BjornerBrenti}: the faces form a regular cell decomposition of the boundary of a simple convex polytope of dimension $k-1$.

@@def:KA@@
For a commutative realization with \emph{distinct} labels $\ell_1<\dots<\ell_k$, labels are addressed by rank, the histories of a piece $P(m_0,Q)$ are the admissible orderings $A=A_Q(m_0)\subseteq S_k$, and $G_A$ is the swap graph of the piece, since swapping adjacent distinct labels is the transposition of adjacent positions. For contents with repeated labels the swap graph is defined but the permutohedral complex is not; Betti numbers are claimed only for distinct-label pieces.

@@def:prefix-defined@@

@@lem:faces-prefix@@

\section{What requires arithmetic}\label{sec:needs}
\Cref{tab:needs} is the contract that the later chapters keep: a result in a row is claimed only for the stated class of systems. In particular the reconstruction, duality and information-ladder results hold for \emph{arbitrary} finite sets of histories, the nerve, block and cube results for every \emph{prefix-defined} system (which contains SELECT-II but is much larger), and the closure, threshold and survival laws use the divisibility lattice of SELECT-II.
@@RAW:ch02.tex::\begin{table}[ht]\centering\small::\end{table}@@

\section{Probability conventions}\label{sec:probconv}
Three different probabilities occur and are never mixed. On a fibre the uniform distribution is used unless another is stated, so that $\Ent(W|T=t)=A(t)$ and every entropy of \cref{ch:ancestry} is an exact finite sum of terms $\frac nN\log_2\frac Nn$. In \cref{ch:separation} a \emph{system} (a content with a source regime) is itself a random variable and one asks how much an observation tells about a property such as $\beta_2>0$; the \emph{natural} ensemble is uniform over the systems of a stated census and the \emph{balanced} ensemble gives equal weight to the two values of the property, both being experimental conventions rather than laws of the grammar. Finally, entropies computed on a whole finite ensemble by exact class counts are \emph{in-sample}, whereas the information a rule learned on one sample gives about another is \emph{out-of-sample}; \cref{sec:predictive} shows that they differ substantially.
